\documentclass[../../main.tex]{subfiles}

\begin{document}
得益于类似“品牌效应”的宣传，很多同学在学习经典力学、量子力学、电动力学之前多少都会了解它们在干什么，即使不了解，简单学习也会理清学科脉络，但热力学与统计物理学（大多时候我们简称热统）却没有这样的关注度，热统自身的抽象性再加上中规中矩的课本和讲课老师，同学们往往会迷失在诡异的推导和庞杂的计算中。

吉林大学物理学院的同学第一次接触热统应该是在大一某名师的课堂上，这估计不是一段很好的回忆。我下面要讲的内容类似当年的内容，不过大家已经有过两年的物理学习，再加上我力图站在同学们的视角展开论述，相信同学们会有一个好的阅读题体验。

在讨论热统在干什么之前，大家回想一下所学的物理理论，发现大多在干两件事：一是阐述对自然界的认识，即我们如何认识世界；二是阐述自然界运动（更准确的说是演化）的规律，即我们如何改造世界。本科物理的核心课程基本都涵盖这两方面，但有一个除外：即热学\footnote{热学在本章语境有两个含义，一个是热学学科，一个是普通物理热学课程，请读者自行区分}和热统。

热学和热统是本科物理中学习热学学科的两门核心课程，两者是密不可分的，下面我从课程内容的角度回答本章的核心问题。

作为一个学科，热学研究宏观系统的热现象。在本科阶段，这个学科重点在阐述我们对热平衡系统的认识，而不是关注系统是怎么从一个热平衡态到另一个热平衡态。热学的核心研究对象是热力学系统和热力学过程，在本科阶段，前者主要是平衡态系统，后者主要是准静态过程，无论从哪方面，我们学习的重点都是描述平衡态。

描述热力学系统\footnote{我们主要关注平衡态,但不代表这些手段只能用在平衡态上}的手段，大家都有所了解，分为两种，一种从宏观着手，称为热力学；另一种从微观着手，称为统计力学。

热力学大家在热学课程中已经接触，可能留给大家的印象就是一团气体一会膨胀一会收缩。但其实热力学是涵盖极广的学问，某种意义上，实际生产生活接触的所有系统，都可以用热力学描述，包括最常见的气体，还有溶液、电磁介质、橡皮筋乃至化学反应体系等各种复杂系统。热力学力图从系统的宏观量入手，总结这些宏观量的关系，热学学科主要聚焦于简单的气体系统；而热统中的热力学着眼任意系统并借助更多的数学手段，也就是热力学被劈开成两部分在两门课中学习。\footnote{具体来说，热力学系统的基本概念、借助理想气体引入的热力学定律、简单的相变描述、对固体液体的简单介绍以及基础的非平衡输运属于热学课程；而基于任意系统的热力学定律、平衡热力学系统的深入数学描述、多元多组分系统及相变构成了热统课程的热力学部分。}

统计力学类似，但划分方式不是内容的深度而是研究视角。统计力学有两种视角：针对微观粒子的统计和针对系统的统计。前者大多数人比较熟悉且觉得理所当然，比如计算系统的内能就要统计各微观粒子的能量；后者就比较反常，它的思路是根据我们研究的系统，假想出一系列相关的系统，系统的内能来自对这些假想的系统的能量的统计。第一种视角发展在统计力学的早期，主要研究气体分子，因此称为气体动理论。这种视角难以处理分子间相互作用，后续被第二个视角取代\footnote{不过，这种思路在现代借助计算机被发扬光大，形成一个研究微观世界的有力手段：分子动力学模拟}。至于第二种视角，假想出的大量系统，合在一起称之为系综，因此这种统计称为系综统计，是现代统计力学的主要内容。系综统计\footnote{不负责任的说，这种统计观点可以认为是把一个系统看成粒子，假想的系统就是无相互作用的粒子}着眼系统，粒子间的相互作用也属于系统的一部分，所以它天然就能解决分子间相互作用的问题。从课程内容来说，分子动理论在热学中学习，而系综统计在热统的统计力学\footnote{国内热统教科书在统计力学部分往往还会重点讲述一种统计观点，即最概然分步法，这种观点是气体动理论的延续，缺陷相同而且其结论可以从系综统计中以更优雅更严格的方式推出，故本讲义只会简略提及}中学习。

至此我们已经回答了“热统在干什么”，而且可能也回答了部分人看到热统的课程名冒出来的一个问题“热学课程已经学习了热力学和统计力学，为什么到了四大力学阶段又把名字并在一起再学一遍”。我从课程内容展开论述，一方面是因为我们通过课程学习物理，课程内容本身就阐明了物理内容的逻辑；另一方面提醒大家注意热统中的热力学绝不是热学的复习。经典力学课程中的力学可以说是普通力学的复习，但如前所述，热统中的热力学比热学中的热力学更深入更数学化，完全可以说是新知识\footnote{有部分同学以为热力学部分就是热学的复习，没认真对待，再加上统计力学繁杂的计算，导致期末考试热力学不会写、统计力学算不完}。

课程的学习是离不开教材的，我推荐刘川老师的《热力学和统计物理学》以及一本外国教材《热物理概念 热力学与统计物理学》。

最后要强调一点，本章和之后的章节构成了一个热统笔记，不过是讲义风格的笔记。它本质上是笔记，会比任课老师所讲的内容多，也会缺少一些任课老师讲的内容，任课老师不会参考本笔记出题，本笔记也不是为了备考而准备，任何人因为阅读本笔记导致的考试丢分本人概不负责。当然，笔记中的错误和读者不理解的地方都可以给我发邮箱\footnote{QQ邮箱：2944052508@qq.com}交流。
\end{document}
